% Add "answers" option to print the answers
\documentclass[addpoints]{exam}

\usepackage[utf8]{inputenc}
\usepackage{calc}
\usepackage[dvipsnames]{xcolor}
\usepackage{ifthen}
\usepackage{tcolorbox}

% \makebox[\textwidth]{Name:\enspace\hrulefill}
%
% \vspace{5mm}
%\makebox[\textwidth]{Instructor’s name:\enspace\hrulefill}

\definecolor{MyLightGray}{gray}{0.60}

\newcommand{\truefalse}
{		\begin{checkboxes}
			 \choice True
			 \choice False. Explanation:
		\end{checkboxes}
		\vspace{10pt}
}


% Check if 'answers' is an option
\newcommand{\ifanswers}[2]{%
  \ifthenelse{\boolean{exam@answers}}{#1}{#2}%
}



\newcommand{\answerspace}[1]
{
	\ifprintanswers
	\else
		\setlength{\mylength}{10pt}%
		% \vspace{#1\mylength}
		% Make the grey box
		\par\nobreak\noindent
		\makebox[\linewidth][l]{%
			% set inner padding to zero, so only the frame shows
			{\setlength{\fboxsep}{0pt}%
			% draw a color frame, white (or transparent) inside
			\fcolorbox{MyLightGray}{white}{%
			% the inner box has zero width content but fixed height
			\parbox[t][#1\mylength][c]{%
				% subtract twice the frame‐thickness from the full width
				\dimexpr\linewidth-2\fboxrule\relax
			}{\strut}%
			}%
			}%
		}%
		\par
		\vspace{3\mylength}

  	\fi
}
\newlength{\mylength}

\newcommand{\code}[1]{{\texttt{#1}}}


\usepackage{color}
\usepackage{lastpage}
\usepackage{verbatim}
\usepackage{fancyvrb}
\usepackage{graphicx}
\usepackage{float}
\usepackage{minted}
\usepackage{multicol}

\RecustomVerbatimCommand{\VerbatimInput}{VerbatimInput}{}

\newcommand{\Code}[1]{\mintinline{python}{#1}}

\renewcommand{\solutiontitle}{\noindent\textbf{Solution sketch: }}

\definecolor{SolutionColor}{rgb}{0.8,0.9,1}

\cfoot{Page \thepage\ of \pageref{LastPage}}

\rhead{Initials: \rule{20mm}{0.4pt}}


%%%%%%%%%%%%%%%%%%%%%%%%%%%%%%%%%%%%%%%%
\begin{document}
\fvset{numbers=left,numbersep=3pt}
\thispagestyle{empty}

\begin{center}
	{\LARGE CS 408 Fall 2026 -- Grey-Box Fuzzing Quiz \\}
\fbox{\fbox{\parbox{5.5in}{\centering
	\begin{itemize}
        \item Try to come up with a solution on your own first, but
          after that talk to your friends and neighbors.

        \item Full points will be rewarded for \textbf{any} (reasonable)
          attempt.
        \end{itemize}
	}}}
\end{center}

\vspace{5mm}
\makebox[\textwidth]{\textbf{Name (please print)}: \rule{100mm}{0.8pt} }

\vspace{10mm}
\makebox[\textwidth]{\textbf{PUID}: \rule{100mm}{0.8pt} \hspace{35mm}}

\begin{questions}

\question The following function crashes when its input contains the
  substring \Code{"fuzz"}. We've labelled each branch for reference:

\begin{minted}{python}
def scan(s: str) -> None:
    matched = 0
    for c in s:
        if c == "fuzz"[matched]:
            matched += 1             # Match: (C4, BB5)
            if matched == 4:
                crash()              # Bad: (C6, BB7)
        else:
            matched = 0              # Reset: (C4, BB9)
\end{minted}

  Suppose that our sole seed input is a random string of 4 characters,
  each of which was uniformly drawn from the set of 95 printable ASCII
  characters.

\begin{parts}

\part What is the likelihood that a \textbf{black-box} fuzzer whose
  sole mutation is swapping on the seed's characters with a randomly
  chosen printable ASCII character triggers the failure? Roughly how
  many attempts would it need to find the bug? %
  \answerspace{6} %
  \begin{solution} %
    All four characters must be right, each with probability $1/95$:
    \[
      P = \frac{1}{95^4} = \frac{1}{81{,}450{,}625}
    \]
    so about \textbf{81 million} attempts on average. Nothing is learned
    from a failure --- \Code{"fuzy"} and \Code{"qqqq"} look equally bad,
    because neither crashed.
  \end{solution}

\part Now suppose the fuzzer records a snapshot of the \textbf{test
    path} that each input induces, in the form of the sequence of
  branches it executes, in order, and adds any newly mutated input to
  the seed set when its path differs from all the other snapshots seen
  so far.

  For each input below, in this order, say whether it is kept or
  discarded. %

  \begin{center}
    \begin{tabular}{llll}
      1. \Code{"x"} \quad & 2. \Code{"f"} \quad & 3. \Code{"fu"} \quad \\[2pt]
      4. \Code{"xf"} \quad & 5. \Code{"fx"} \quad & 6. \Code{"fuz"} \\
    \end{tabular}
  \end{center}
  \answerspace{8} %
  \begin{solution} %
    \textbf{Every single one is kept.} The paths are all different:

    \begin{center}\small
      \begin{tabular}{ll}
        \Code{"x"}   & Reset \\
        \Code{"f"}   & Match \\
        \Code{"fu"}  & Match, Match \\
        \Code{"xf"}  & Reset, Match \\
        \Code{"fx"}  & Match, Reset \\
        \Code{"fuz"} & Match, Match, Match \\
      \end{tabular}
    \end{center}

    Note what happened with \Code{"xf"} and \Code{"fx"}: they take the
    same branches the same number of times, differing \emph{only} in
    the order, and neither gets further into \Code{"fuzz"} than
    \Code{"f"} already did. The criterion keeps them both anyway.
  \end{solution}

%\newpage

\part Recall that AFL does not record the test path. Instead it
  records, for each branch (i.e., edge out of a conditional node in
  the CFG), \textbf{how many times} that branch ran, bucketed into
  ranges, and keeps an input when some branch lands in a bucket it has
  not seen before. How does that change your answer to (b)? %
  \answerspace{20} %
  \begin{solution} %
    First, the counts. Only the three labelled edges are recorded, and
    \Code{Bad} is never taken, since none of these inputs is a complete
    \Code{"fuzz"}:

    \begin{center}\small
      \begin{tabular}{lccc}
        & Match & Bad & Reset \\ \hline
        \Code{"x"}   & 0 & 0 & 1 \\
        \Code{"f"}   & 1 & 0 & 0 \\
        \Code{"fu"}  & 2 & 0 & 0 \\
        \Code{"xf"}  & 1 & 0 & 1 \\
        \Code{"fx"}  & 1 & 0 & 1 \\
        \Code{"fuz"} & 3 & 0 & 0 \\
      \end{tabular}
    \end{center}

    What happens next depends on exactly what ``a bucket it has not
    seen'' is compared against. Both readings are reasonable, and both
    were accepted.

    \medskip
    \textbf{Reading 1 --- compare the whole vector of counts.} An input
    is kept when its triple of counts differs from every triple seen so
    far. Then \Code{"xf"} is \textbf{kept}: $(1,0,1)$ is a combination
    neither \Code{"x"} $(0,0,1)$ nor \Code{"f"} $(1,0,0)$ produced.
    \Code{"fx"} is then \textbf{discarded}, because its counts are
    identical to \Code{"xf"}'s --- the two differ only in the order of
    the branches, and order is precisely what counting throws away.

    \medskip
    \textbf{Reading 2 --- compare each (edge, bucket) pair separately},
    which is what AFL actually does: it keeps one global table of every
    pair it has ever seen, and an input survives only if it contributes
    a pair that is not already in it. Then \textbf{both} \Code{"xf"}
    and \Code{"fx"} are discarded: \Code{Match} at 1 came from
    \Code{"f"}, \Code{Reset} at 1 came from \Code{"x"}, and neither
    input adds anything beyond those two.

    \medskip
    Either way, the ladder survives: \Code{Match} running 1, 2 and 3
    times falls in three different buckets, so \Code{"f"},
    \Code{"fu"} and \Code{"fuz"} are kept under both readings. And
    either way the fuzzer discards at least one input that the test-path
    criterion in (b) kept --- an input that made no progress at all.
  \end{solution}

\part Why do you think AFL uses bucketed hit counts rather than the
  full test path? %
  \answerspace{20} %
  \begin{solution} %
    Because the set of test paths is effectively unbounded, while the
    set of hit-count signatures is not.

    Every distinct length and every distinct ordering of branches is a
    new path, so ``keep anything with a new path'' keeps almost
    \emph{everything}: a 500-byte input has a 500-iteration path that
    nothing else will match. The queue grows without limit, and most of
    what is in it --- like \Code{"xf"} and \Code{"fx"} --- represents no
    progress at all. The fuzzer would spend its budget re-mutating
    noise.

    Hit counts throw away the ordering and squash the magnitudes into
    eight buckets. That keeps the signal (looping once, twice or three
    times usually \emph{does} mean the input got structurally further)
    and discards the noise (94 versus 95 iterations almost never does).
    The result is a bounded signature that can be stored as a fixed-size
    bitmap and compared in constant time, which is what makes it
    affordable to do on every single execution.

    This is the same trade as everywhere else in this course: a coarser
    criterion is cheaper and less precise. Here the imprecision is the
    whole point.
  \end{solution}

\end{parts}

\end{questions}

\end{document}
